% =====================================================================
% CHAPTER 6: Results and Discussion (RESTRUCTURE_PLAN_V2, department
% template dated 2026-07-14). RED: 6.1 Introduction, 6.last Summary.
% BLACK middle headings chosen from the template's suggested set
% (Experimental Results, Comparison, Error Analysis, Validation,
% Performance Evaluation, Discussion), populated only where old
% mainmatter/ch8.tex ("Results and Discussion" under the prior
% 9-chapter restructure) has matching text:
%   - Investigation and Analysis of Data / Training Dynamics
%     -> 6.2 Experimental Results (subsection)
%   - Investigation and Analysis of Data / In-domain SceneFlow numbers
%     -> 6.2 Experimental Results (subsection)
%   - SceneFlow + Middlebury 2014 comparison tables and discussion
%     -> 6.3 Comparison
%   - Rectification tolerance, qualitative inspection, 3D reconstruction
%     -> 6.4 Error Analysis
%   - Design-space / architecture / training-recipe ablation studies
%     -> 6.5 Validation (these validate the design choices of Chapter
%     \ref{ch:design}; the department calls this slot "Validation")
%   - Efficiency and deployment latency/memory numbers
%     -> 6.6 Performance Evaluation
%   - Old "Discussion" section -> 6.7 Discussion, unchanged in substance
% Dataset descriptions are NOT repeated here; they live in Chapter
% \ref{ch:experimental} (Section \ref{sec:exp_datasets}) and are only
% pointed to. Every \label from the old chapter is preserved verbatim;
% see the SOURCES/DROPPED/MOVED-AWAY block at the end of this file.
% =====================================================================
\chapter{Results and Discussion}
\label{ch:results}

\section{Introduction}
This chapter reports the accuracy, efficiency, and cross-domain
behaviour of the trained model. We cover four threads: in-domain
accuracy on FlyingThings3D against published methods, zero-shot
accuracy on Middlebury 2014 against matched-protocol checkpoints,
on-device latency against the real-time target, and the training
dynamics and ablation studies behind the architectural design of
Chapter \ref{ch:design}. Unless stated otherwise, the results come
from a single production run trained on the full SceneFlow
FlyingThings3D corpus, under the protocol and configuration of Table
\ref{tab:train_config} in Chapter \ref{ch:implementation}. We follow
the accuracy metrics of \eqref{eq:metrics}; latency is measured at
batch one, 384 by 640 resolution; and cross-domain behaviour is
assessed on Middlebury 2014, which the model never saw during
training. Chapter \ref{ch:experimental} (Section
\ref{sec:exp_datasets}) sets out both datasets, the physical stereo
rig, and the evaluation protocol in full; this chapter does not repeat
those details.

\FloatBarrier
\section{Experimental Results}

\subsection{Training Dynamics}
Figure \ref{fig:training_curves} shows the training trajectory. Loss
falls sharply during the OneCycle warm-up, then descends steadily as
the learning rate anneals toward its floor. Validation end-point error
tracks the loss closely, bottoming out at 0.79 pixels near step 53,000
and holding flat between 0.80 and 0.81 pixels for several thousand
steps as the rate nears zero. We took the best checkpoint at step
53,000, not the final step 60,000, since further training gave no
improvement. Figure \ref{fig:convergence} shows the same pattern: broad
surfaces settle early while object boundaries and thin structures keep
sharpening. The run survived several cloud preemptions: scheduler and
optimizer state are checkpointed and restored exactly, so a preemption
costs at most the steps since the last validation. The one-cycle
schedule's high peak rate acts as a regularizer mid-training, visible
as the sustained loss variance in Figure \ref{fig:training_curves}; its
long cosine descent then turns that exploration into steady gains,
with more than half the final improvement coming after the rate drops
below a quarter of its peak. The per-batch loss wobble in the
high-rate phase is the schedule working as designed, not instability,
and the flat validation tail confirms that 60,000 steps was about
right for this model and dataset.

\begin{figure}[H]
\centering
\includegraphics[width=\textwidth]{fig_4_1_training_curves}
\caption{Training loss and validation accuracy curves over the training run.}
\label{fig:training_curves}
\end{figure}

\begin{figure}[H]
\centering
\includegraphics[width=\textwidth]{fig_4_2_convergence}
\caption{Predicted disparity across training checkpoints for two FlyingThings3D scenes.}
\label{fig:convergence}
\end{figure}

\FloatBarrier
\subsection{In-Domain Accuracy: SceneFlow Test}
We evaluated the trained model on the full FlyingThings3D test set,
4,370 pairs at native 960 by 540 resolution, using the best checkpoint.
Table \ref{tab:sceneflow_results} reports 0.78 pixel end-point error,
3.40 percent D1-all, and a median error of 0.13 pixels; four frames
with no valid ground-truth pixel after the disparity cutoff are
excluded from the macro-average. The metric spread tells its own
story: the median of 0.13 pixels is six times smaller than the mean of
0.78, and the RMSE of 3.64 is nearly five times larger, so most pixels
are highly accurate while a small population of large outliers sets
the average. The bad-rate ladder agrees: errors above half a pixel
affect 15.3 percent of pixels, but above three pixels only 4.0
percent, so most failures are moderate and catastrophic ones rare.
Figure \ref{fig:sceneflow_qualitative} shows where those outliers
live: occlusion boundaries and thin structures, consistent with
recurrent reasoning that runs at quarter resolution and an upsampler
that can recombine disparity but cannot invent it.

\begin{table}[!htbp]
\centering
\caption{Full FlyingThings3D test-set results across all eight evaluation metrics.}
\label{tab:sceneflow_results}
\small
\begin{tabular}{cccccccc}
\toprule
EPE (px) & RMSE & Median & bad-0.5 (\%) & bad-1 (\%) & bad-2 (\%) & bad-3 (\%) & D1-all (\%)\\
\midrule
0.781 & 3.64 & 0.130 & 15.32 & 8.92 & 5.34 & 4.00 & 3.40\\
\bottomrule
\end{tabular}
\end{table}

\begin{figure}[H]
\centering
\includegraphics[width=\textwidth]{fig_4_3_sceneflow_qualitative}
\caption{Left image, ground truth, prediction, and error map for held-out FlyingThings3D scenes.}
\label{fig:sceneflow_qualitative}
\end{figure}

\FloatBarrier
\section{Comparison}

\subsection{Comparison with Published Methods on SceneFlow}
Table \ref{tab:sceneflow_comparison} places our numbers alongside
published methods, though the comparison is not strictly
protocol-identical: each published EPE follows its own paper's
SceneFlow protocol, while our value comes from the full test set at
native resolution. The model falls between the efficient and iterative
families: more accurate than the classical 3D-cost-volume baseline
PSMNet, close to the efficient LightStereo-S, and a fraction of the
parameters of the iterative and foundation-scale systems.

\enlargethispage{3cm}
\vspace{-0.5cm}
\begin{table}[H]
\centering
\caption{Comparison with published SceneFlow methods (protocol not directly comparable to ours).}
\label{tab:sceneflow_comparison}
\small
\renewcommand{\arraystretch}{0.9}
\begin{tabular}{lcccc}
\toprule
Method & Params (M) & SF EPE (px) & Latency (ms) & Hardware\\
\midrule
PSMNet \cite{chang2018psmnet} & 5.2 & 1.09 & 410 & Titan Xp\\
HITNet (L) \cite{tankovich2021hitnet} & 0.97 & 0.43 & 54 & Titan V\\
CoEx \cite{bangunharcana2021coex} & 2.72 & 0.69 & 27 & RTX 2080Ti\\
RAFT-Stereo \cite{lipson2021raftstereo} & 11.23 & 0.61 & 380 & RTX 6000\\
IGEV-Stereo \cite{xu2023igev} & 12.60 & 0.47 & 180 & RTX 3090\\
LightStereo-S \cite{guo2025lightstereo} & 3.44 & 0.73 & 17 & RTX 3090\\
FoundationStereo \cite{wen2025foundationstereo} & ${\sim}340$ & 0.34 & ${\sim}470$ & RTX 4090\\
\midrule
StereoLite (ours) & 2.96 & 0.78 & 49.8 (fp16) & RTX 3050 laptop\\
\bottomrule
\end{tabular}
\end{table}

\FloatBarrier
\newpage
\subsection{Zero-Shot Cross-Domain: Middlebury 2014}
Cross-domain accuracy is the real test for an edge model: one accurate
only on its training domain has little field use. We evaluated the
trained model zero-shot on the 23 perfect-set scenes of Middlebury
2014, unseen during training, under a protocol identical for every
method in Table \ref{tab:mb14_comparison}, reaching 1.71 pixel
end-point error and 10.9 percent D1-all, the smallest model in the
comparison at 2.96 million parameters. A gap remains to LiteAnyStereo
(6.9 percent) and sixteen-iteration IGEV-Stereo (5.0 percent), both
larger models; LiteAnyStereo also adds a three-stage
knowledge-distillation pipeline that our supervised-only recipe
skips. Figure \ref{fig:mb14_perscene} shows sixteen of the 23 scenes
within the reference band, with the aggregate dragged up by a few hard
scenes; Table \ref{tab:mb14_perscene_full} lists the per-scene split.

\begin{table}[!htbp]
\centering
\caption{Zero-shot Middlebury 2014 comparison across three models (our reproduction of public checkpoints; T4 fp32 eager latency).}
\label{tab:mb14_comparison}
\small
\resizebox{\textwidth}{!}{%
\begin{tabular}{lcccccccc}
\toprule
Model & Params (M) & T4 (ms) & EPE & bad-0.5 & bad-1 & bad-2 & bad-3 & D1-all (\%)\\
\midrule
\textbf{StereoLite (ours)} & 2.96 & 86 & \textbf{1.71} & 39.6 & 23.9 & 14.5 & 11.2 & \textbf{10.9}\\
LiteAnyStereo \cite{jing2025liteanystereo} & 7.60 & 64 & 1.17 & 35.1 & 19.7 & 10.5 & 7.4 & 6.9\\
IGEV-Stereo (16 it.) \cite{xu2023igev} & 12.60 & 305 & 0.86 & 22.2 & 12.4 & 7.2 & 5.2 & 5.0\\
\bottomrule
\end{tabular}}
\end{table}

\begin{figure}[H]
\centering
\includegraphics[width=\textwidth]{fig_4_4_mb14_perscene}
\caption{Per-scene zero-shot D1-all on Middlebury 2014, sorted by error.}
\label{fig:mb14_perscene}
\end{figure}

The easiest scenes are ordinary furnished rooms: Adirondack at 2.4
percent D1-all, Recycle at 2.9, and Storage at 4.2, comparable to the
reference models. The hardest four alone raise the mean by several
points: Jadeplant at 34.9 percent, Flowers at 30.3, Mask at 24.7, and
Pipes at 17.9. These scenes share dense thin repeated structures such
as leaves and rods, plus true disparities near or beyond the
192-pixel evaluation cap, exactly where a quarter-resolution,
bounded-search prediction performs worst. Figure
\ref{fig:mb14_qualitative} shows both ends of that spectrum. The
model's cross-domain failure is not diffuse: it is localized to an
identifiable scene type, which a future distillation direction could
target directly.

\begin{figure}[H]
\centering
\includegraphics[width=0.8\textwidth]{fig_4_5_mb14_qualitative}
\caption{Left image, ground truth, and prediction for the easiest and hardest Middlebury 2014 scenes.}
\label{fig:mb14_qualitative}
\end{figure}

\FloatBarrier
\section{Error Analysis}

\subsection{Robustness to Rectification Imperfection}
\label{sec:rectification}
A second thesis objective is tolerance to imperfect rectification. We
evaluated the trained model on a 400-pair test subset with the right
image shifted vertically by a controlled offset and ground truth
unchanged, emulating residual vertical misalignment between views.
Table \ref{tab:rectification} and Figure \ref{fig:rectification}
report the result. The zero-offset baseline of 1.03 pixels is measured
on this 400-pair perturbation subset rather than the full 4,370-pair
test set, so it sits above the 0.78-pixel headline of Table
\ref{tab:sceneflow_results}; the quantity of interest here is the rise
under misalignment, not the absolute level. Accuracy degrades
gracefully across the practical range. At a half-pixel offset,
end-point error rises only from 1.03 to 1.22 pixels and D1-all from
4.3 to 4.6 percent; at a full pixel, error stays at 1.53 pixels with
6.2 percent D1-all. A calibrated rig typically holds residual vertical
error below one pixel, so accuracy survives this range; the tolerance
comes from the refinement stage's local correlation lookup, which
checks a small neighborhood around each match instead of assuming a
horizontal epipolar line. Beyond this range degradation is steep: a
two-pixel offset raises D1-all to 15.8 percent, a four-pixel offset to
41.2 percent, since the network lacks explicit vertical search and
large misalignment pushes matches outside the lookup window. This is
measured tolerance to small misalignment, not immunity to arbitrary
de-rectification.

\begin{table}[!htbp]
\centering
\caption{Accuracy on FlyingThings3D under simulated vertical rectification offset.}
\label{tab:rectification}
\small
\begin{tabular}{lcccc}
\toprule
Vertical offset (px) & EPE (px) & bad-1 (\%) & bad-2 (\%) & D1-all (\%)\\
\midrule
0.0 & 1.03 & 10.28 & 6.63 & 4.29\\
0.5 & 1.22 & 14.62 & 7.58 & 4.60\\
1.0 & 1.53 & 26.27 & 11.84 & 6.19\\
2.0 & 2.43 & 47.73 & 27.71 & 15.83\\
4.0 & 5.07 & 71.59 & 55.68 & 41.17\\
\bottomrule
\end{tabular}
\end{table}

\begin{figure}[!htbp]
\centering
\includegraphics[width=0.72\textwidth]{fig_4_9_rectification}
\caption{End-point error and D1-all against vertical rectification offset.}
\label{fig:rectification}
\end{figure}

The synthetic sweep is controlled with known ground truth, but the
condition it emulates occurs for real on a low-cost rig. To confirm
this on physical data, we ran the model on 997 rectified pairs from
the project's stereo rig, a commodity camera whose factory
rectification leaves residual vertical misalignment, and compared it
against the FoundationStereo teacher's pseudo-ground-truth on the same
pairs. Agreement is close: mean end-point error is 1.45 pixels and
D1-all is 11.1 percent, matching the in-domain synthetic error despite
the real rig's imperfect rectification. Since both see the same
misaligned input, this figure measures agreement with a strong
reference rather than absolute error, corroborating rather than
replacing the controlled sweep. Together, the two results confirm that
the model withstands the rectification imperfections of an
inexpensive camera.

\FloatBarrier
\subsection{Qualitative Results}
Visual inspection reinforces the quantitative results above. On the
synthetic test set, Figure \ref{fig:sceneflow_qualitative} (shown
earlier in this chapter) shows accurate reconstruction of broad
surfaces, with error confined to occlusion boundaries. On the unseen
Middlebury scenes, Figure \ref{fig:mb14_qualitative} shows clean
reconstruction of easy scenes, with failure localized to thin repeated
structures rather than spread across the image. The plane-rendering
step of Chapter \ref{ch:design} shows at the boundaries: predicted
disparity follows object edges more closely than a bilinear upsample
would. Figure \ref{fig:camera} extends this inspection to real data:
four indoor pairs from the calibrated stereo rig, none used in
training or fine-tuning, yield plausible disparity maps that cleanly
separate foreground from background, and the model transfers to a
real camera without adaptation.

\begin{figure}[!htbp]
\centering
\includegraphics[width=\textwidth]{fig_4_6_camera}
\caption{Left image and predicted disparity for real indoor stereo pairs.}
\label{fig:camera}
\end{figure}

\FloatBarrier
\subsection{3D Reconstruction of Real Scenes}
This subsection presents the promised reconstruction. We
back-projected disparity from three real indoor scenes to metric point
clouds via the calibrated relations of \eqref{eq:backprojection},
cleaned them with a statistical outlier filter and voxel downsampling
per Section \ref{sec:reconstruction_method}, and rendered them from
three viewpoints (Figure \ref{fig:reconstruction}). The reconstructions
are metrically plausible. In the first scene, two fire extinguishers
stand proud of the wall and door in the oblique and top-down views. In
the second, the corridor recovers a smooth depth gradient to the
bright far end. In the third, the metal cabinet sits in front of the
back wall at its correct standoff. The flying points along depth
boundaries are the three-dimensional footprint of the same
occlusion-edge error visible in the error maps of Figure
\ref{fig:sceneflow_qualitative}; the depth quantization on distant
surfaces reflects depth's inverse relation to disparity, since a fixed
sub-pixel disparity step spans a growing depth interval with distance.
These pairs come from the same low-cost rig whose imperfect
rectification we examined in Section \ref{sec:rectification}; we did
no fine-tuning on the rig's data, so the clouds are the zero-shot
output of the checkpoint evaluated throughout this chapter.

\begin{figure}[H]
\centering
\includegraphics[width=\textwidth]{fig_4_10_reconstruction}
\caption{Left image and point clouds from three viewpoints for real indoor scenes.}
\label{fig:reconstruction}
\end{figure}

\FloatBarrier
\subsection{Comparison with LiDAR Point Clouds}
\label{sec:lidar_comparison}
It is worth placing these reconstructions next to the sensor they are
usually measured against. LiDAR is the accuracy standard for depth: it
measures range directly by timing a laser pulse, and its points are
metrically exact to within a centimetre or two. But that accuracy comes
one point at a time. A scanning unit samples the scene along a fixed
handful of rotating beams, so its returns land on a small number of scan
rings and the space between them is simply not measured.
Figure~\ref{fig:lidar_compare} shows this directly: a raw Velodyne scan
from KITTI \cite{geiger2012kitti}, on the left, is accurate everywhere it
has a point, yet any single surface in it carries only a thin scatter of
returns; our stereo cloud of the fire-extinguisher scene, on the right,
covers the same kind of surface continuously, giving the full shape of
each object and the layout of the scene rather than a sample of it.

For a mobile robot, that difference in density is not a cosmetic one. A
robot planning a path or building a map has to know the shape and extent
of the obstacles around it, not just the range to a few isolated points
on them, and a thin object such as a table edge, a chair leg, or the very
extinguisher nozzles in Figure~\ref{fig:reconstruction} can fall between a
LiDAR's scan lines entirely while a dense stereo surface still resolves
it. Density, in other words, buys spatial understanding, and it is
exactly what a passive stereo cloud provides where a low-cost LiDAR does
not.

Against this, the accuracy gap is a reasonable price. The reconstructions
here are not centimetre-exact the way a LiDAR return is; their error
grows with distance, as the inverse disparity-to-depth relation of
Section~\ref{sec:reconstruction_method} requires, and their depth
boundaries carry the flying points visible in the figure. What the trade
buys is substantial: a passive stereo rig costs a small fraction of a
scanning LiDAR, draws a few watts instead of tens, emits nothing into the
scene, and returns a dense surface instead of a sparse one, all from two
commodity image sensors. For the edge robots, drones, and assistive
devices this thesis targets, where budget, power, and the shape of nearby
obstacles matter more than millimetre-accurate ranging, trading a bounded
loss of point accuracy for dense, low-cost, passive coverage is the right
exchange, and it is the one the whole design is built around.

\begin{figure}[!htbp]
\centering
\includegraphics[width=\textwidth]{fig_6_lidar_compare}
\caption[Sparse LiDAR point cloud versus our dense stereo point cloud.]{A real sparse LiDAR scan (KITTI \cite{geiger2012kitti}, left) beside our dense stereo reconstruction of the fire-extinguisher scene (right). LiDAR is accurate but samples a surface along a few scan rings; the stereo cloud covers it densely, trading a bounded accuracy loss for dense, low-cost, passive coverage.}
\label{fig:lidar_compare}
\end{figure}

\FloatBarrier
\section{Validation}
This section validates the design decisions behind the architecture of
Chapter \ref{ch:design} through a sequence of controlled ablation
studies, run under the matched-harness protocol described there,
rather than in a single design pass. We group the studies in three:
early exploration that selected the encoder, refinement scheme, and
loss on the earlier tile chassis; the ablation of the recurrent family
that produced the final model; and training-recipe and efficiency
studies on the final model.

\subsection{Design-Space Exploration of the Tile Chassis}
An earlier campaign on a smaller tile chassis settled three questions.
Table \ref{tab:ablation_encoder} compared a GhostConv stack against two
truncated detection backbones under a matched 20-pair harness; YOLO26s
won on every metric, 0.528 pixel end-point error against 0.625 for
GhostConv and 0.712 for the smaller YOLO26n, at a two-millisecond
latency cost, and was selected for the final model. Table
\ref{tab:ablation_refinement} tested the iterated coarse-to-fine
scheme against two alternatives: extending refinement to finer scales
gained almost nothing (0.573 against 0.578 pixels) at 37 percent more
latency, while HITNet-style single-pass-per-scale propagation
collapsed accuracy to 0.844 pixels. Iterated refinement at coarse
scales is thus the load-bearing choice, kept in the final model. Table
\ref{tab:ablation_loss} compared nine formulations under the same
harness; the threshold-stack objective with the D1-relative hinge won
on end-point error, bad-2, and D1, and is the ancestor of the
production objective of \eqref{eq:loss}. The clearest negative result:
sequence-weighted multi-scale L1, which imitates the recurrent
literature's iteration-weighted loss without exposing true
per-iteration outputs, was the worst of all nine at 1.005 pixels.
Faking a sequence loss across scales does not work; the final model
instead exposes true multi-scale outputs to the objective.

\begin{table}[!htbp]
\centering
\caption{Encoder ablation on the tile chassis under the matched 20-pair harness.}
\label{tab:ablation_encoder}
\small
\begin{tabular}{lccccc}
\toprule
Encoder & Params (M) & EPE & bad-1 (\%) & D1-all (\%) & Latency (ms)\\
\midrule
GhostConv stack & 0.538 & 0.625 & 13.31 & 3.54 & 23.5\\
YOLO26n (truncated) & 0.808 & 0.712 & 18.07 & 2.90 & 24.1\\
\textbf{YOLO26s (truncated)} & 2.061 & \textbf{0.528} & \textbf{10.88} & \textbf{2.50} & 25.8\\
\bottomrule
\end{tabular}
\end{table}

\begin{table}[!htbp]
\centering
\caption{Refinement-scheme ablation on the tile chassis under the matched harness.}
\label{tab:ablation_refinement}
\small
\begin{tabular}{lccccc}
\toprule
Scheme & Params (M) & EPE & bad-1 (\%) & D1-all (\%) & Latency (ms)\\
\midrule
\textbf{iterated coarse-to-fine} & 0.538 & \textbf{0.578} & \textbf{10.72} & 3.44 & \textbf{23.8}\\
+ refinement at 1/2 & 0.576 & 0.573 & 11.12 & \textbf{3.25} & 32.5\\
single pass per scale & 0.487 & 0.844 & 19.30 & 4.79 & 25.2\\
\bottomrule
\end{tabular}
\end{table}

\vspace{-0.5cm}
\begin{table}[H]
\centering
\caption{Loss-formulation ablation across nine objectives on the tile chassis.}
\label{tab:ablation_loss}
\small
\begin{tabular}{lcccc}
\toprule
Objective & EPE & bad-0.5 (\%) & bad-2 (\%) & D1-all (\%)\\
\midrule
plain L1 & 0.663 & 27.6 & 6.09 & 3.99\\
sequence-weighted multi-scale L1 & 1.005 & 53.5 & 13.25 & 4.04\\
L1 + gradient & 0.618 & \textbf{23.9} & 5.93 & 3.92\\
L1 + bad-1 hinge & 0.763 & 45.8 & 6.21 & 3.65\\
cocktail (L1 + grad + bad-1) & 0.692 & 40.0 & 6.12 & 3.68\\
cocktail, 0.5-px hinge & 0.676 & 38.3 & 6.21 & 3.50\\
threshold stack & 0.672 & 38.0 & 5.97 & 3.45\\
\textbf{threshold stack + D1 hinge} & \textbf{0.591} & 31.5 & \textbf{5.64} & \textbf{3.35}\\
Charbonnier & 0.631 & 31.0 & 5.78 & 3.58\\
\bottomrule
\end{tabular}
\end{table}

\FloatBarrier
\subsection{Architecture Ablation of the Recurrent Family}
The second campaign built the recurrent context family that became the
final model, changing one mechanism at a time on a matched 20-pair
harness (Table \ref{tab:ablation_gev4line}). The base design reached
0.424 pixel end-point error; a learned update gate plus the YOLO26s
encoder brought it to 0.319. Rejected additions: a selective recurrent
unit (0.474), a fresh quarter-resolution initialization (0.455), and
the sharp-tail and high-resolution-refine variants, which gave minor
per-metric gains at added latency without improving the headline
error. The decisive addition, a quarter-resolution geometry encoding
volume with fail-soft fusion, reached 0.261 pixels, the campaign best,
at 0.041 million added parameters. Step and batch budgets differ
slightly across runs, so the table is an engineering record, not a
controlled comparison; the next subsection's matched studies
re-validate these conclusions.

\begin{table}[!htbp]
\centering
\caption{Architecture ablation of the recurrent family that produced the final model.}
\label{tab:ablation_gev4line}
\small
\begin{tabular}{lccccc}
\toprule
Variant & Params (M) & EPE & Median & bad-1 (\%) & Latency (ms)\\
\midrule
recurrent context (base) & 2.831 & 0.424 & 0.167 & 7.62 & 10.7\\
+ geometry-guided state & 2.919 & 0.459 & 0.218 & 8.16 & 36.7\\
+ update gate & 2.921 & 0.473 & 0.246 & 7.75 & 12.2\\
gate, YOLO26n encoder & 1.410 & 0.387 & 0.182 & 5.98 & 10.9\\
gate, YOLO26s encoder & 2.921 & 0.319 & 0.127 & 5.24 & 11.9\\
+ sharp tail & 2.950 & 0.322 & 0.136 & \textbf{5.21} & 14.6\\
+ high-resolution refine & 3.106 & 0.351 & 0.186 & 5.18 & 16.1\\
+ selective recurrent unit & 3.398 & 0.474 & 0.255 & 8.20 & 14.8\\
+ raw 1/4 initialization & 2.948 & 0.455 & 0.195 & 10.64 & 20.7\\
\textbf{+ GEV and fail-soft fusion} & 2.962 & \textbf{0.261} & \textbf{0.059} & 5.25 & 24.6\\
\bottomrule
\end{tabular}
\end{table}

\FloatBarrier
\subsection{Training-Recipe and Efficiency Studies}
Four further controlled studies (Figure \ref{fig:ablations}) justify
the training recipe and efficiency pass of the final model. The
training recipe, not any architectural block, was the largest accuracy
lever: in a matched nine-arm comparison, the OpenStereo augmentation
triplet lowered the best validation end-point error from 2.778 to
1.921 pixels, a 31 percent reduction well above the noise floor. The
edge-inference optimization pass (Table \ref{tab:ablation_efficiency})
held accuracy while cutting latency by roughly a third:
exact-equivalence fixes plus the narrow geometry volume brought
latency from 46.6 to 30.2 milliseconds, with end-point error unchanged
within noise (2.908 to 2.906 pixels). An exact-equivalence class
(zero-copy views, batched lookups, hoisted static channels, fused
heads) is bitwise or mathematically identical, so its accuracy deltas
are pure noise; the narrow geometry band, the one accuracy-relevant
change, was validated separately as neutral. The boundary-sharpness
study (Table \ref{tab:ablation_blurfix}) selected plane rendering,
which cut the bad-0.5 rate from 55.36 to 42.91 percent, the sharpest
and fastest result of the study, since rendering from planes replaces
a costlier path. Adding the bimodal head or top-k initialization on top
of plane rendering worsened results (44.40 and 46.58 percent bad-0.5):
the fixes anti-synergize, so we adopted only plane rendering, a
negative result that constrains future work. Table
\ref{tab:ablation_native} reports the input-protocol study on both
axes: training on native crops cut the native-axis end-point error
from 6.672 to 2.869 pixels versus the legacy global-resize protocol,
at no cost on the resized axis, so we adopted the native-crop protocol
for the production run.

\begin{table}[!htbp]
\centering
\caption{Inference-efficiency ablation on the 100-pair harness, 9k steps (RTX 3050 fp32: 106.7 to 61.4 ms, Table \ref{tab:efficiency_findings}).}
\label{tab:ablation_efficiency}
\small
\begin{tabular}{lccccc}
\toprule
Variant & Params (M) & EPE & bad-1 (\%) & D1-all (\%) & Latency (ms)\\
\midrule
base (gev4) & 2.962 & 2.908 & 30.32 & 16.31 & 46.6\\
+ exact-equivalence fixes & 2.962 & 3.021 & 32.09 & 17.79 & 35.8\\
+ narrow GEV band & 2.962 & 2.906 & 32.29 & 16.88 & 30.2\\
\bottomrule
\end{tabular}
\end{table}

\begin{table}[!htbp]
\centering
\caption{Boundary-sharpness fixes and their composition on a 500-pair windowed split.}
\label{tab:ablation_blurfix}
\small
\begin{tabular}{lcccc}
\toprule
Variant & EPE & bad-0.5 (\%) & D1-all (\%) & Latency (ms)\\
\midrule
control (narrow) & 2.006 & 55.36 & 12.57 & 22.2\\
+ bundle1 (top-k init) & 1.982 & 49.95 & 12.14 & 21.9\\
+ bimodal aux head & 1.779 & 43.26 & 11.53 & 25.2\\
\textbf{+ plane rendering (adopted)} & 1.826 & \textbf{42.91} & 12.19 & 15.0\\
\midrule
plane + bimodal & 1.847 & 44.40 & 11.62 & 25.5\\
plane + bimodal + bundle1 & 1.853 & 46.58 & 11.93 & 17.2\\
\bottomrule
\end{tabular}
\end{table}

\begin{figure}[H]
\centering
\includegraphics[width=\textwidth]{fig_4_7_ablations}
\caption{Ablation results for augmentation, inference optimization, boundary blur fixes, and input protocol.}
\label{fig:ablations}
\end{figure}

\begin{table}[!htbp]
\centering
\caption{Input-protocol ablation on the dual-axis evaluation (50 validation frames).}
\label{tab:ablation_native}
\small
\begin{tabular}{lcccccc}
\toprule
& \multicolumn{3}{c}{Native axis ($960\times 540$)} & \multicolumn{3}{c}{Resized axis ($640\times 384$)}\\
\cmidrule(lr){2-4}\cmidrule(lr){5-7}
Arm & EPE & bad-1 & D1-all & EPE & bad-1 & D1-all\\
\midrule
global resize (legacy) & 6.672 & 58.69 & 35.81 & 1.783 & 29.91 & 11.44\\
\textbf{native crop (adopted)} & \textbf{2.869} & \textbf{36.75} & \textbf{16.47} & 1.783 & \textbf{29.30} & \textbf{10.71}\\
native full frame & 2.967 & 38.68 & 16.74 & 2.557 & 36.45 & 16.36\\
\bottomrule
\end{tabular}
\end{table}

\FloatBarrier
\section{Performance Evaluation}

\subsection{Efficiency and Deployment Measurements}
Table \ref{tab:latency} reports the measured efficiency. On the RTX
3050 laptop GPU at batch one and 384 by 640 resolution, the optimized
model ran in 49.8 milliseconds at half precision and 61.4 milliseconds
at full precision, with inference working sets of 0.26 and 0.35
gigabytes. On the target device, a Jetson Orin Nano, the optimized
eight-bit model measured 36.3 milliseconds per frame on-device, or
27.5 frames per second at the same resolution. This nears the 30
frames per second real-time target, confirming interactive rates on
the intended edge accelerator, not just the laptop GPU. The on-device
working set stays well within the platform's memory budget.

\begin{table}[!htbp]
\centering
\caption{Latency and memory across RTX 3050 and Jetson Orin Nano configurations (Jetson memory is estimated).}
\label{tab:latency}
\small
\begin{tabular}{lcccc}
\toprule
Configuration & Latency (ms) & FPS & Inference memory (GB)\\
\midrule
RTX 3050, fp32 & 61.4 & 16.3 & 0.35\\
RTX 3050, fp16 & 49.8 & 20.1 & 0.26\\
Jetson Orin Nano, INT8 TensorRT & 36.3 & 27.5 & ${\sim}0.15$\\
\bottomrule
\end{tabular}
\end{table}

\FloatBarrier
\section{Discussion}
These results answer the four research questions posed in Chapter
\ref{ch:intro}. The model reaches 0.78 pixel end-point error within a
few million parameters, answering the first. It generalizes zero-shot
to an unseen domain at 10.9 percent D1-all, short of the
distillation-trained references, answering the second. It runs
interactively on a laptop GPU and reached 27.5 frames per second on
the Jetson Orin Nano, near the real-time target, answering the third.
The fourth question, tolerance to imperfect rectification, is
addressed by the perturbation study of Section \ref{sec:rectification}.
Figure \ref{fig:pareto_ours} places the model on the zero-shot
accuracy-efficiency front, at the low-parameter end below the larger
distillation-trained and iterative references. This comparison bears
on an open question about cross-domain generalization: LiteAnyStereo,
the strongest efficient reference, is itself non-iterative and owes
its accuracy to a three-stage distillation recipe rather than to
iteration, which suggests that the training pipeline carries more of
the cross-domain gain than architectural iteration does, though our
study does not isolate the two factors directly. A few limitations
bound these conclusions: the results come from a single run, so
run-to-run variance is unknown; the model was not submitted to the
KITTI leaderboard; and it lacks left-right consistency or occlusion
reasoning. These limits are taken up as future directions in the
concluding chapter.

\begin{figure}[!htbp]
\centering
\includegraphics[width=0.78\textwidth]{fig_4_8_pareto_ours}
\caption{Middlebury 2014 D1-all against trainable parameters for three zero-shot models.}
\label{fig:pareto_ours}
\end{figure}

\FloatBarrier
\section{Summary}
This chapter evaluated the model on every objective axis. In domain,
it reached 0.78 pixel end-point error and 3.40 percent D1-all on the
FlyingThings3D test set. Zero-shot on unseen Middlebury 2014 scenes, it
reached 10.9 percent D1-all under a protocol identical to the
reference models, the smallest network in that comparison; remaining
failures cluster in scenes with thin repeated structures and
disparities near the evaluation cap. Under controlled vertical
rectification error, end-point error rose from 1.03 to 1.53 pixels
across the sub-pixel to one-pixel range that a calibrated rig
produces, and the real-rig agreement study confirmed that tolerance on
physical data. On the Jetson Orin Nano, the optimized eight-bit model
reached 36.3 milliseconds per frame, near the 30 frames per second
real-time goal; real-camera disparity maps and metric point-cloud
reconstructions confirmed that the checkpoint transfers to physical
data without adaptation.

% SOURCES: mainmatter/ch8.tex (entire file, all sections: Introduction;
%   Investigation and Analysis of Data [Training Dynamics,
%   Design-Space Exploration of the Tile Chassis, Architecture
%   Ablation of the Recurrent Family, Training-Recipe and Efficiency
%   Studies]; Results [In-Domain Accuracy: SceneFlow Test, Zero-Shot
%   Cross-Domain: Middlebury 2014, Robustness to Rectification
%   Imperfection, Efficiency and Deployment Measurements, Qualitative
%   Results, 3D Reconstruction of Real Scenes]; Discussion; Summary).
%   All 22 old \label anchors (10 figures, 11 tables/subsection label,
%   1 section label sec:rectification) plus the new chapter label
%   ch:results are preserved verbatim; none renamed, none dropped.
% DROPPED: nothing. Every subsection of the old chapter has a home
%   under the new black headings.
% MOVED-AWAY: nothing content-wise; only headings were reorganized.
%   The old "Investigation and Analysis of Data" intro paragraph was
%   split: the training-dynamics framing stayed with Section 6.2
%   (Experimental Results), and the ablation framing was rewritten as
%   the new Section 6.5 (Validation) intro, since training dynamics
%   and the ablation campaigns no longer share one heading. The
%   SceneFlow comparison table and its "falls between the efficient
%   and iterative families" sentence were split out of the old
%   "In-Domain Accuracy" subsection into their own Comparison
%   subsection; the raw in-domain numbers (Table tab:sceneflow_results,
%   the metric-spread and bad-rate discussion, Figure
%   fig:sceneflow_qualitative) stayed under Experimental Results.
%   Cross-references to chapter numbers were resolved to labels that
%   now exist elsewhere in mainmatter2/: "Chapter 3" (architecture and
%   matched harness) -> \ref{ch:design}; "Chapter 7" (train config
%   table) -> \ref{ch:implementation} via Table \ref{tab:train_config};
%   "the four research questions from Chapter 3" -> \ref{ch:intro},
%   since the ownership map moves problem statements/research
%   questions to Chapter 1; "Chapter 9" (future directions,
%   distillation) -> \ref{ch:conclusion} (the concluding chapter's
%   label was not yet defined by its owning agent at write time;
%   \ref{ch:conclusion} follows this book's ch:intro/ch:litreview/
%   ch:design/ch:implementation/ch:experimental/ch:results naming
%   convention and should be verified against the actual label chosen
%   in the new ch9.tex during the B4 QA pass). Dataset descriptions
%   were not reproduced; the Introduction points to Chapter
%   \ref{ch:experimental}, Section \ref{sec:exp_datasets} instead.
